% ============================================================================
% Chapter 12: Storage Devices
% ============================================================================

\chapter{Storage Devices}
\label{ch:storage}

\begin{objectives}
  \item Name and describe different types of storage devices
  \item Compare storage devices by capacity, speed, and portability
  \item Understand the difference between primary and secondary storage
  \item Explain what cloud storage is and how it works
\end{objectives}

\bytetip[tip]{Storage is where your computer keeps everything --- from your homework to your favourite photos. Think of it like a digital wardrobe for all your stuff!}

\section{Types of Storage}
\label{sec:storage-types}
\index{storage!types}

\begin{theorybox}[Primary vs Secondary Storage]
\begin{itemize}[leftmargin=*, label={\textcolor{ModuleBlue}{\faAngleRight}}, itemsep=4pt]
  \item \textbf{Primary storage}\index{primary storage} (RAM) --- Temporary, fast, loses data when powered off
  \item \textbf{Secondary storage}\index{secondary storage} (HDD, SSD, USB) --- Permanent, keeps data even when off
\end{itemize}
\end{theorybox}

% --- Figure 12.1: Storage Capacity Bar Chart ---
\begin{figure}[H]
\centering
\begin{tikzpicture}
\begin{axis}[
  xbar, bar width=14pt,
  xlabel={Typical Capacity (GB)},
  xlabel style={font=\small\sffamily},
  symbolic y coords={DVD, USB 4GB, USB 64GB, SSD 512GB, HDD 1TB, Cloud},
  ytick=data,
  yticklabel style={font=\small\sffamily},
  xticklabel style={font=\small\sffamily},
  nodes near coords,
  nodes near coords style={font=\tiny\sffamily\bfseries},
  xmin=0, xmax=1100,
  width=12cm, height=7cm,
  enlarge y limits=0.15,
]
\addplot[fill=FunCoral!70, draw=FunCoral] coordinates {(4.7,DVD)};
\addplot[fill=FunSky!70, draw=FunSky] coordinates {(4,USB 4GB)};
\addplot[fill=FunMint!70, draw=FunMint] coordinates {(64,USB 64GB)};
\addplot[fill=FunPurple!70, draw=FunPurple] coordinates {(512,SSD 512GB)};
\addplot[fill=ModuleBlue!70, draw=ModuleBlue] coordinates {(1000,HDD 1TB)};
\addplot[fill=LabGreen!70, draw=LabGreen] coordinates {(1000,Cloud)};
\end{axis}
\end{tikzpicture}
\caption{Comparing storage device capacities (Cloud shown as 1000+ GB)}
\label{fig:storage-chart}
\end{figure}

% --- Table 12.1: Storage Device Comparison ---
\begin{table}[H]
\centering
\caption{Comparing different storage devices}
\label{tab:storage-compare}
\renewcommand{\arraystretch}{1.3}
\begin{tabularx}{\textwidth}{l l X l l}
\toprule
\rowcolor{HeaderRow}
\textcolor{white}{\textbf{Device}} & \textcolor{white}{\textbf{Type}} & \textcolor{white}{\textbf{Capacity}} & \textcolor{white}{\textbf{Speed}} & \textcolor{white}{\textbf{Portable}} \\
\midrule
\rowcolor{RowLight} HDD & Magnetic & 500 GB -- 4 TB & Medium & \textcolor{WarnRed}{\ding{55}} \\
\rowcolor{RowDark} SSD & Flash & 128 GB -- 2 TB & Fast & \textcolor{WarnRed}{\ding{55}} \\
\rowcolor{RowLight} USB Drive & Flash & 4 GB -- 256 GB & Medium & \textcolor{LabGreen}{\ding{51}} \\
\rowcolor{RowDark} DVD & Optical & 4.7 GB & Slow & \textcolor{LabGreen}{\ding{51}} \\
\rowcolor{RowLight} Cloud & Online & Unlimited & Varies & \textcolor{LabGreen}{\ding{51}} \\
\bottomrule
\end{tabularx}
\end{table}

\section{Cloud Storage}
\label{sec:cloud}
\index{cloud storage}

\hl{Cloud storage} lets you save files on the internet instead of your computer. Services like \textbf{Google Drive}, \textbf{OneDrive}, and \textbf{Dropbox} let you access your files from any device, anywhere!

\begin{realworld}
When you save a photo to Google Photos or upload a document to Google Drive, you're using cloud storage. The file isn't on your phone or laptop --- it's on a powerful computer (called a \textbf{server}) in a data centre, possibly thousands of kilometres away!
\end{realworld}

\bytetip[fun]{If you stacked all the data stored in the world's cloud servers as printed paper, the pile would reach from Earth to Pluto --- and back --- \textbf{multiple times}!}

\begin{funfact}
The first USB flash drive was sold in 2000 and could store just \textbf{8 MB} of data --- that's not even enough for one high-quality photo today! Modern USB drives can hold up to \textbf{2 TB}, which is 250,000 times more.
\end{funfact}

\begin{reviewqs}
  \item What is the difference between primary and secondary storage?
  \item Name four types of secondary storage devices.
  \item Which storage device is the fastest: HDD, SSD, or DVD?
  \item What is cloud storage, and name two cloud storage services.
  \item Why is an SSD faster than an HDD?
\end{reviewqs}

\begin{challenge}
Check how much free storage your school computer has. Right-click on \textbf{This PC} $\rightarrow$ \textbf{Properties} (or check the C: drive). Write down: (a) Total capacity, (b) Used space, (c) Free space. Calculate what percentage is used!
\end{challenge}
